% SYNC-ID: introduction
\section{引言}\label{sec:introduction}
% CONTENT_DEFERRED_TO_WRITING_PHASE；可见文字仅为非实质性工程占位。
\textit{(THIS IS PLACEHOLDER TEXT, PLZ FULFILL IT WHEN WRITING.)}
% SYNC-ID: introduction-placeholder-1
开头段落将在后续写作阶段说明经确认的背景与范围。本占位段使用普通论文文字，以便检查最终文档的版式，而不是检查一组宏的外观。
% SYNC-ID: introduction-placeholder-2
下一段预留简短动机和写作阶段选定的来源指针；此处不表达任何特定领域的事实。
% SYNC-ID: introduction-placeholder-3
本占位章节的结尾段预留经批准的范围边界，并与中文伴随版保持对应。
% SYNC-ID: introduction-placeholder-4
相邻段落还预留从背景过渡到论文提纲的位置。该过渡应当简洁、可追溯到批准的计划，并在计划审核前不包含任何论断。
% SYNC-ID: introduction-placeholder-5
为了检查版式占位，引言还预留了过渡句、范围限定语和指向后续章节的指针。正式措辞应说明哪些内容属于范围、哪些内容有意排除，以及每个表述在哪里可以核查。这些都是写作任务，不是研究发现，因此本段不包含经验或理论内容。
% SYNC-ID: placeholder-figure
% 占位图源文件 -- 待办：替换为作者批准的语言无关图。
\begin{figure}[t]
  \centering
  \fbox{\parbox[c][1.05in][c]{0.86\columnwidth}{\centering\textit{占位图 -- 待办：在后续写作阶段插入经确认的总览图。}}}
  \caption{占位图 -- 待办：冻结图的事实来源后再替换为正式图注。}
  \label{fig:placeholder-overview}
\end{figure}

